\chapter[Depth and Lefschetz theorems]
{Depth and Lefschetz theorems in~\'etale cohomology} \label{XIV}
\begin{center}
by Mme M. Raynaud\sfootnote{After unpublished notes of A.~Grothendieck.}
\end{center}

\vspace*{1cm}
In \numero \Ref{XIV.1}, we define a notion of ``\emph{\'etale depth}'' which is the analogue in \'etale cohomology of the notion of depth studied in \Ref{III}, in the cohomology of coherent sheaves. After a technical part, we prove in \numero \Ref{XIV.4} some ``Lefschetz theorems'', the central theorem being \Ref{XIV.4.2}. Let $X$ be a scheme, $Y$ a closed subset of $X$, $U$ the complementary open set $X-Y$ and $F$ an abelian sheaf on $X$, for the \'etale topology; in a general way the aim of the Lefschetz theorems is to show that, if $F$ satisfies certain local conditions on $U$, expressible in terms of \'etale depth at the points of $U$, then, under certain additional conditions of a global nature on $U$ (for example $U$ \emph{affine}), the natural map of \'etale cohomology groups
$$
\H^i(X, F)\to \H^i(Y, F_{|Y})$$
is an isomorphism for values $i<n$, where $n$ is a certain integer which is made explicit. Taking for $F$ a constant sheaf, one thus obtains conditions for $\pi_0(X)$ to be equal to $\pi_0(Y)$ and conditions for the abelianized fundamental groups of $X$ and $Y$ to be the same. In \numero \Ref{XIV.5}, the introduction of a notion of ``geometric depth'' allows one to give useful special cases of the Lefschetz theorems (\Ref{XIV.5.7}). Finally in \numero \Ref{XIV.6}, we mention some conjectures, concerning in particular ``noncommutative'' variants of the theorems obtained.
